\documentclass{article}
\usepackage[T1]{fontenc}
\usepackage{lmodern}
\usepackage{graphicx}
\usepackage{amsmath,amssymb}
\usepackage[a4paper,margin=2cm]{geometry}
\usepackage[absolute,overlay]{textpos}
\setlength{\TPHorizModule}{1mm}
\setlength{\TPVertModule}{1mm}
\usepackage[indonesian]{babel}
\usepackage{hyperref}
\setlength{\emergencystretch}{2em}
% unit: Halaman 22

\begin{document}

% === Halaman 22 (llm) ===

\begin{itemize}
    \item OTP Tidak cocok untuk komunikasi dinamis. Bayangkan untuk komunikasi modern seperti:
    \begin{itemize}
        \item WhatsApp
        \item HTTPS
        \item VPN
        \item cloud storage
        \item email
        \item komunikasi IoT, dll
    \end{itemize}
    data yang dikirim terus berubah, perangkat bisa berganti, pengguna bisa \textit{offline}, dan pesan memiliki ukuran yang berbeda-beda.
    \vspace{5mm}
    \item Dengan OTP, setiap tambahan data membutuhkan tambahan kunci yang benar-benar acak dan belum pernah digunakan. Ini membuat manajemen kunci menjadi sangat berat.
\end{itemize}

\end{document}
